% Introduction for UAI - condensed

Knowledge graph (KG) embeddings enable efficient reasoning over structured knowledge~\citep{bordes2013translating,trouillon2016complex}.
As KGs evolve---new entities emerge, existing entities appear in novel relationships---robust uncertainty quantification becomes critical.

\paragraph{The Relation-Agnostic Limitation.}
Recent probabilistic methods learn distributions over embeddings~\citep{chen2019embedding,chen2021probabilistic,Chen2021PERM}, where rare entities have higher variance.
But these variances are \emph{relation-agnostic}: entity $e$ receives a single variance $\sigma^2_e$ regardless of relational context.
Consider $e$ appearing in 1000 triples with relations $\{r_1, \ldots, r_{10}\}$ but never with $r_{11}$.
The model assigns low variance (well-observed), yet queries $(?, r_{11}, e)$ involve genuinely uncertain predictions.

This causes systematic failures (Figure~\ref{fig:main_results}): score-based methods achieve 0.99 AUROC detecting random corruptions (implausible triples) but only 0.44--0.64 on temporal OOD (emerging entities and novel relational contexts that simulate temporal distribution shift).

\paragraph{Two OOD Types.}
We identify two distinct scenarios:
(1) \textbf{Emerging entities}: New/rare entities with poorly-constrained embeddings---high variance correctly signals uncertainty.
(2) \textbf{Novel contexts}: Well-observed entities in unobserved relations---variance is low (falsely confident), but the entity-relation pair was never trained.

\paragraph{Contributions.}
We identify a systematic limitation in probabilistic KG embeddings: learned variances are relation-agnostic despite training on data containing entity-relation co-occurrence patterns.
This failure reveals a fundamental mismatch between standard training objectives (link prediction accuracy) and OOD detection requirements (capturing structural observation patterns).

Our contributions are:
\textbf{(1) Problem identification and impossibility result:} We identify a systematic blind spot in probabilistic KG embeddings---the inability to detect novel relational contexts---and prove that \emph{any} relation-agnostic uncertainty estimator achieves near-random detection on this category (Theorem~\ref{thm:impossibility}).
\textbf{(2) Uncertainty decomposition framework:} We formalize complementary semantic and structural uncertainty signals and show neither alone suffices (Theorem~\ref{thm:complementarity}). This decomposition is the core contribution; the combination mechanism (CAGP) is intentionally minimal.
\textbf{(3) Empirical validation:} We demonstrate this limitation persists across existing methods (0.44--0.64 AUROC on temporal shift despite 0.99 on random corruptions), and that the decomposition consistently achieves 0.94--1.00 AUROC across three benchmarks.

\begin{figure}[t]
\centering
\includegraphics[width=0.85\linewidth]{figures/fig1_main_results.pdf}

\caption{Temporal OOD detection (AUROC). Existing probabilistic methods achieve only 0.44--0.64; CAGP combines complementary semantic and structural signals to reach 0.94--1.00 across three benchmarks.}
\label{fig:main_results}
\end{figure}
